\chapter{Conclusion}

This thesis investigated binary radiographic pneumonia detection in an emergency department--linked cohort derived from MIMIC-CXR-JPG and MIMIC-IV-ED. The study compared triage-only clinical baselines, an image-only DenseNet121 classifier, and an early-fusion multimodal model combining triage features with chest radiographs. A central design principle was the explicit anchoring of prediction at the time of imaging, with all inputs restricted to information available at or before that moment. This constraint ensured that the task corresponds to a well-defined and clinically meaningful decision point and that reported results are aligned with a realistic deployment-time information setting.

On a shared patient-level temporal test partition, the image-only model achieved the strongest overall performance across both discrimination and threshold-based metrics. In particular, the image-only model achieved AUROC 0.746 and AUPRC 0.724, compared with 0.736 and 0.714 for the multimodal model. Paired patient-level bootstrap comparisons of multimodal minus image produced uncertainty intervals for both AUROC and AUPRC that included zero, indicating no reliable advantage from fusion in this setting. Both deep learning models substantially outperformed the clinical baselines, demonstrating that chest radiographs provide significantly stronger predictive signal for this task than the selected triage variables.

Beyond these primary results, the thesis provides additional insight into the structure of the problem itself. Performance was shown to depend strongly on the composition of the negative class, with a marked decrease when distinguishing pneumonia from other abnormal radiographs compared with relatively normal cases. These subgroup effects indicate that the benchmark is intrinsically ambiguous: separating pneumonia from other abnormal thoracic findings is substantially harder than separating pneumonia from relatively normal studies. Calibration analysis further demonstrated that discrimination and probabilistic reliability are distinct properties, with the multimodal model exhibiting somewhat improved calibration under fixed-bin metrics despite lower ranking performance. Decision curve analysis confirmed that both image-based models provide higher utility than the clinical baselines across a range of thresholds, while again showing no clear overall advantage of multimodal fusion over the image-only system.

Taken together, these findings challenge the assumption that multimodal integration is inherently beneficial. In the present setting, where tabular inputs are limited to pre-imaging triage data and strict temporal constraints are enforced, a strong image model captures most of the useful predictive signal. Multimodal fusion, at least in its current form, introduces additional complexity without yielding a consistent performance gain. This result does not imply that multimodal learning is ineffective in general, but rather that its value depends critically on the nature, timing, and complementarity of the available modalities.

More broadly, this work highlights the importance of rigorous problem formulation and evaluation design in medical machine learning. Explicit temporal alignment, leakage-aware feature selection, and patient-level evaluation are not merely implementation details; they fundamentally shape both the difficulty of the task and the validity of the conclusions. Under the strict temporal constraints used here, methodological validity determined which conclusions could be trusted: robust image-only superiority over clinical baselines, but no reliable multimodal gain.

In summary, this thesis demonstrates that careful control of data, time, and evaluation is as important as model architecture in determining the outcome of medical AI systems. The results emphasize that strong unimodal baselines should be established before claiming multimodal benefit, that pneumonia detection in this setting is intrinsically ambiguity-rich, and that empirical validation under realistic constraints is essential for drawing reliable conclusions. These findings should be understood as internal to the present MIMIC-based benchmark, but they provide a strong foundation for future work on alternative targets, richer structured inputs, and broader external validation.